% \begin{savequote}
% \begin{verbatim}
% ( Excellent time to become a missing )
% ( person.                            )
%  ------------------------------------
%   o
%    o
%        ___  
%      {~._.~}
%       ( Y )
%      ()~*~()   
%      (_)-(_) 
% \end{verbatim}
% \end{savequote}
\def\ChapterCode{TOX}
\chapter
[\CO{[TOX]} Vergiftungen]
{Vergiftungen}
\ChapterIntro
{%
~~~
}
{%
\begin{description}
  \item[Maintainer:] Sebastian Gabriel
  \item[Autoren:] Diverse
  \item[Reviewer:] Standard-Reviewprozess
  \item[Version:] {\VarDocVersionTypeName} ({\VarDocVersionTypeDescription})
  \item[SHA1:] {\DocGitCommit}
\end{description}

{\TxTeilkapitelAASS}
}
% %%%%%%%%%%%%%%%%%%%%%%%%%%%%%%%%%%%%%%%%%%%%%%%%%%%%%%%%%%%
% \section{Vergiftungen -- Toxikologie}
\label{topic:vergiftung}



\clearpage

\Dictum[Hohenheim von Theoprast]{Die Dosis macht das Gift.}

\DefinitionInPara{Eine \T{Vergiftung} ist eine (meist  dosisabhängige)
organische  Funktionsbeeinträchtigung bzw. Schädigung des
Organismus durch eine aufgenommene Substanz.}
Vergiftungen sind eine interessante Sache: Ein an sich
nützlicher Stoff kann sich in eine todbringende Waffe
verwandeln, wenn er im Übermaß zugeführt wird. So ist zum
Beispiel ein Überleben ohne Salz nicht möglich. Führt man
jedoch 20\,dag zu, endet das meist tödlich. Sonst nützliche
Medikamente können schon bei einer Dosisabweichung im
Mikrogramm-Bereich eine fatale Wirkung haben.

Andererseits können Vergiftung in Abhängigkeit vom
ursächlichen Stoff so ziemlich alle Erscheinungsformen
darbieten. Oft findet man nur \EE{unspezifische Symptome}
vor, wie z.\,B. Erbrechen, Durchfall, Schmerzen im
Magen-Darm-Trakt, Bewusstseinsstörungen, Krampfanfälle oder
Entgleisung der Vitalparameter.
Es ist weder möglich, noch sinnvoll, auf sämtliche mögliche
Vergiftungen im Detail einzugehen. Im Folgenden werden
einige häufig anzutreffende Vergiftungen sowie
Herangehensweisen für andere, weniger häufige, vorgestellt
und besprochen. 





\Text{Der Weg in den Körper -- Aufnahme des Giftes}{%
Die Aufnahme der Substanz kann über verschiedene Wege
erfolgen:
\begin{itemize}
  \item \EE{Oral}: über den Mund
  \item \EE{Inhalativ}: Einatmen
  \item \EE{Stechen, spritzen}: in die Vene (intravenös,
  z.\,B.
  Drogenkonsum), in den Muskel, in die %(Unter-)
  Haut (z.\,B. Insektenstich)
  \item \EE{Über die Haut}: transkutan; fettlösliche
  Substanz kann  Hautbarriere durchdringen (z.\,B.:  organische
  Lösungsmittel,  Insektenschutzmittel etc.)
\end{itemize}
}{}{}{}{}





% \subparagraph*{Vorgehensweise}
% 
% 5-Fingerregel:
% 
% \begin{enumerate}
%   \item Vitalparameter
%   \item Giftentfernung
%   \item Antidot-Therapie (Notarzt)
%   \item Gift-Asservierung
%   \item Transport
% \end{enumerate}

\TextSlide{Vergiftungsinformationszentrale}{%
Bei der Vergiftungsinformationszentrale kann Expertenrat
eingeholt werden, sofern der verantwortliche Stoff bekannt
ist. Unter Umständen kann wird hier auch bei der
Identifizierung des Stoffes geholfen.
 }{
\begin{itemize}
  \item Expertenrat
\end{itemize}

}{}{}{}

\opt{REG-AUT}{\Fazit{Vergiftungsinformationszentrale: 01 /
406 43 43}}




Gegen die meisten Vergiftungen kann man präklinisch,
ursächlich wenig unternehmen. Das Hauptaugenmerk liegt
daher auf: 


% M %%%%%%%%%%%%%%%%%%%%%%%%%%%%%%%%%%%%%%%%%%%%%%%%%%% M %
% Vergiftungen, allgemein
\ProcedurePrintDefault{MT65091C}


\clearpage % TEMP temp clearpage
\section{Vergiftungen mit Alkohol}

\Layout{20}{Einleitung}{%
Die Vergiftung mit Alkohol ist häufig anzutreffen, weil sie
sozial akzeptiert ist.
Wenn der Betroffene \enquote{über's Ziel hinaus schießt},
kann es schnell zu lebensbedrohlichen Zuständen kommen.
}
{}
{./images/TAE/dif-w2-zelt-1.jpg}
{}
{}


\Z{Die akute Alkoholvergiftung ist von der chronischen
Alkoholabhängigkeit (Alkoholkrankheit) zu unterscheiden.}

\TextSlide{Wirkung}
{%
Alkohol  wirkt enthemmend, z.\,T. stimulierend und
stimmungsmodulatorisch. Weiters kommt es zu einer Sprach- 
und Gangunsicherheit. 

\bigskip

\bigskip

}{%
\begin{itemize}
    \item Enthemmend, stimmungsmodulatorisch
    \item Sprach- und Gangunsicherheit
    \item Bewusstseinstrübung
\end{itemize}
}{}{}{}

\TextSlide{\TxGefahren}
{%
Mit zunehmender Menge kommt es zu einem totalen
Kontrollverlust, Bewusstseinstrübung,
Schmerzunempfindlichkeit, Temperaturregulationsstörung und
Inkontinenz\footnote{\emph{Inkontinenz:} Unfähigkeit Harn
und/oder Stuhl zu halten.}. Im Extremfall kann es zu einer 
vegetativen Entgleisung und einem tief komatösen Zustand 
kommen. Spätestens dann besteht akute Lebensgefahr. Weitere 
Gefahren sind Atemdepression, Aspiration, Ausfall der 
Schutzreflexe und -- besonders in der kalten Jahreszeit --
Unterkühlung.
}{%
\begin{itemize}
    \item Koma
    \item Atemdepression
    \item  Ausfall der Schutzreflexe
    \item Aspiration
    \item Schmerzunempfindlichkeit
    \item Unterkühlung
\end{itemize}
}{}{}{}


Es ist unbedingt zu erheben, ob noch andere Substanzen
(Medikamente, Drogen, \ldots) eingenommen wurden
(\E{Mischintoxikation}). Aufgrund der oft reduzierten
Schmerzempfindlichkeit ist besonders auf Verletzungen zu
achten.
Die \EE{vitale Bedrohung} ist vor allem anhand der
Vitalparameter, der körperlichen Untersuchung und des
Bewusstseinszustandes zu beurteilen. 

Der Alkoholentzug und das Entzugsdelir wird unter
\FullPrettyRef{topic:alkohol-entzug}, abgehandelt.


% M %%%%%%%%%%%%%%%%%%%%%%%%%%%%%%%%%%%%%%%%%%%%%%%%%%% M %
% Vergiftungen mit Alkohol
\ProcedurePrintDefault{MT51001C}


\Text{Chronische Schäden}
{
\begin{itemize}
    \item Leberschäden, Zirrhose
    \item Aszites (Wasserbauch)
    \item Krampfadern in der unteren Speiseröhre  (Blutungsgefahr)
    \item Verminderte Blutgerinnung
    \item Vitaminmangel
    \item Geistiger Abbau, Gedächtnisschwäche
    \Z{(Korsakow-Syndrom, Wernicke-Enzephalopathie)}
    \item Anämie
    \item Pankreatitis
    \item Fetales Alkoholsyndrom (Schädigung des Fötus in der Schwangerschaft)
    \cite{Jones1973,Guerri2009}
\end{itemize}
}{	

}{}{}{}













\section{Vergiftungen durch Medikamente}

\Layout{61}{}
{%
Überdosierungen von Medikamenten können komplett
unterschiedliche Symptome erzeugen, je nach eingenommenen
Medikamenten. Manchmal können die Symptome sogar erst Stunden
nach der Einnahme auftreten und tödlich enden. Vergiftungen
mit Medikamenten gehören daher immer abgeklärt und die
Patienten hospitalisiert. 

Neben unabsichtlichen Fehleinnahmen oder Überdosierungen
werden Arzneimittel oft auch in Selbstmordabsicht
eingenommen. Die Umstände der Einnahme müssen daher
abgeklärt werden.
}{

}{./images/teaser/imgp0482_00800.jpg}{}{Gabriel}

\Z{
\paragraph{Paracetamol (Mexalen\texttrademark,
APA\texttrademark, \ldots)} Paracetamol kann in hoher Dosis
zu Leberversagen führen. Symptome treten erst nach Stunden
auf. Es gibt ein Gegenmittel, dieses muss jedoch
rechtzeitig gegeben werden. Ohne Behandlung kommt es zum
Leberversagen und in Folge zum Tod. 

\paragraph{Digitalis} Digitalis ist ein auf das Herz
wirkendes Medikament und wird seit langer Zeit zur
Behandlung der Herzinsuffizienz und von
Herzrhythmusstörungen eingesetzt. Es ist jedoch sehr
schwierig die richtige Dosis für einen Patienten zu finden,
daher kommt es -- trotz korrekter Einnahme -- häufig zu
Überdosierungen. Die Folge sind Herzrhythmusstörungen und
andere EKG-Veränderungen und \enquote{Farbensehen}.
}

\section{Vergiftungen mit Suchtmitteln und
\enquote*{Drogen}}

\subsection{Opiate: Heroin \& Co.}

Heroin und Morphium gehören zu der Gruppe der Opiate und
sind im Drogenmilieu weit verbreitet, die Anwendung erfolgt
i.\,d.\,R. intravenös. 

\TextSlide{\TxSymptome{} der Opiatvergiftung}
{%
Die gefährlichste Wirkung der Opiate ist die
Atemdepressionen, d.\,h. die Atemfrequenz wird dosisabhängig
immer weiter reduziert bis es schließlich zu einem
Atemstillstand kommt. Die Patienten zeigen
dementsprechend Zeichen einer Hypoxie
(Zyanose). Charakteristisch ist die starke Verengung der
Pupillen \E{(stecknadelkopfgroße
Pupillen)}\index{Pupillen!stecknadelkopfgroße}. Oft
findet man alte Einstichstellen in der Armbeuge. Da in der
\enquote{Szene} oft unsauberes oder getauschtes Spritzbesteck
verwendet wird, besteht eine erhöhte Erkrankungsrate an HIV
und Hepatitis C. Eine weitere Gefahr geht vom Spritzbesteck
aus, wenn es nicht sicher verwahrt wird (\Ca{Vorsicht vor
der Nadel!})
}{
\begin{itemize}
    \item Atemdepression
    \item Hypoxie
    \item Stecknadelkopfgroße Pupillen
    \item Oft Begleiterkrankungen (HIV, Hep. C) 
    \item \Ca{Vorsicht Spritzbesteck!}
    \item Gegenmittel verfügbar, wirkt aber nicht so lang
\end{itemize}
}{}{}{}

\Text{Gegenmittel}
{%
Jedes Notarztmittel führt ein Opiat-Gegenmittel mit sich,
dieses ist sehr effektiv.\footnote{\Z{Naloxon. Im Handel
als: \emph{Narcanti}\texttrademark,
\emph{Naloxon-ratiopharm}\texttrademark, u.\,a.}} Es ist
jedoch wichtig zu wissen, dass dieses Gegenmittel nicht so
lange wirkt wie das Opiat, d.\,h. es kann wieder zu einer
verminderten Atmung kommen, wenn die Wirkung des
Gegenmittels nachlässt. Daher soll der Transport nur mit
Arztbegleitung erfolgen.

}{

}{}{}{}


% M %%%%%%%%%%%%%%%%%%%%%%%%%%%%%%%%%%%%%%%%%%%%%%%%%%% M %
% Vergiftung mit Opiaten
\ProcedurePrintDefault{MT40021C}


\subsection{\enquote{Partydrogen}: Ecstacy, Schwammerl,
Special K \ldots}

\TextSlide{\TxBeschreibung}
{Als \T{Partydrogen} bezeichnet man Amphetaminpräparate oder
ähnliche Stimulantien. Zu den bekanntesten Drogen dieser Art
gehören  \I{Ecstasy} (\I{XTC}, \I{E}, \I{X}\Z{; MDMA
(3,4-Methylen\-dioxy\-meth\-amphet\-amin)}). Verwendet wird
die Substanz häufig als \EE{Partydroge}: Der Patient erlebt
ein gesteigertes Glücksgefühl, ein Gefühl von
unerschöpflicher Energie und ein gesteigertes Erleben der
Umwelt.

}{
\begin{itemize}
  \item Ecstasy (XTC, E, X).
  \item Gesteigertes \enquote{Glücksgefühl}.
  \item Gefühl von unerschöpflicher Energie.
  \item Gesteigertes Erleben der Umwelt.
\end{itemize}
}{}{}{}

\TextSlide{\TxSymptome}
{% Die Nebenwirkungen sind jedoch gravierend und können
lebensgefährlich sein. Man kann sie unterteilen in
körperliche und psychische Nebenwirkungen.
Psychisch kommt es dosisabhängig zu Halluzinationen und
Panik-Attacken.
Körperlich können Patienten \E{tachykarde
Herzrhythmusstörungen} zeigen, \E{Krampfanfälle},
\E{hypertensive Krisen}, massiver
\E{Körpertemperaturanstieg} über 40\textdegree{} und
\E{Muskelzerfall}.
Bei einer ausgeprägten Vergiftung kann eine vitale Bedrohung
bestehen. Häufig kommt es zum
\I[Flüssigkeitsmangel!Ecstacy]{Flüssigkeitsmangel}
(\I[Dehydratation!Ecstacy]{Dehydratation}) da der
Patient einerseits auf das Trinken vergisst, und
andererseits i.\,d.\,R. körperlich aktiv ist und dabei durch
Schwitzen viel Flüssigkeit verliert.

\vspace{1.5cm}

}{
\begin{itemize}
  \item Nebenwirkungen Psychisch
  \begin{itemize}
    \item Halluzinationen
    \item Panik-Attacken
  \end{itemize}
  \item Nebenwirkungen Körperlich
  \begin{itemize}
    \item Tachykarde Herzrhythmusstörungen
    \item Hypertensive Krisen
    \item Krampfanfälle
    \item Massiver Temperaturanstieg über
    {\textgreater}\,40\DegC*
    \item Flüssigkeitsmangel
  \end{itemize}
\end{itemize}
}{}{}{}


% M %%%%%%%%%%%%%%%%%%%%%%%%%%%%%%%%%%%%%%%%%%%%%%%%%%% M %
% Vergiftung mit Uppers
\ProcedurePrintDefault{MT43091C}


%„Liquid Ecstasy”: Gamma-Butyrolacton-Entzugsdelir mit
%Rhabdomyolyse und dialysepflichtiger Niereninsuffizienz
\cite{Supady2009}








\clearpage % TEMP temp clearpage
\section{Gase}

\subsection{Vergiftungen mit Stickgasen}

% M %%%%%%%%%%%%%%%%%%%%%%%%%%%%%%%%%%%%%%%%%%%%%%%%%%% M %
% Vergiftung mit Stickgasen
\ProcedurePrintDefault{MT59XX1C}


\subsubsection{Kohlenmonoxid (\texorpdfstring{\ce{CO}}{CO}) -- oder: Von defekten
Heizlüftern, Gasthermen und Autoabgasen}

\Layout{60}{Einleitung}{%
Kohlenmonoxid (\ce{CO}) ist ein farb- und geruchloses Gas.
Esentsteht bei unvollständiger Verbrennung (\ce{O2}-Zufuhr
${\downarrow}$), dies passiert vor allem bei schlecht
gewarteten Gasheizungen und Durchlauferhitzern. \ce{CO}
bindet 300 mal besser an den roten Blutfarbstoff Hämoglobin
als \ce{O2}, daher verdrängen schon kleine Mengen den
Sauerstoff aus dem Blut! Die Folge ist logischerweise eine
Unterversorgung mit Sauerstoff.
}
{}
{./images/gabriel-sebastian-aass/Atemschutz_IMG_2043_00800.jpg}
{}
{GABS}




\TextSlide{Anamnese und Szene}
{
\begin{itemize}
    \item Mehrere, unerklärliche Opfer, leblose Haustiere
    \item Gasofen, Durchlauferhitzer
    \item \emph{\enquote{Immer beim Duschen!}}
\end{itemize}
}{
\begin{itemize}
    \item Mehrere, unerklärliche Opfer, tote Tiere
    \item Gasofen, Durchlauferhitzer
    \item \emph{\enquote{Immer beim duschen!}}
\end{itemize}
}{}{}{}

\TextSlide{\TxSymptome}
{% 
Kohlenmonoxid-beladenes Blut ist -- wie
Sauerstoff-beladenes -- hellrot. Der Patient ist daher eher
nicht bleich oder zyanotisch, sondern die Haut ist oft rosig.
Pulsoxymeter älterer Bauart verwechseln
Kohlenmonoxid-beladenes und Sauerstoff-beladenes Blut, daher
wird ein falsch hoher Sauerstoffsättigungswert angezeigt!

Es stehen die Symptome des \EE{Sauerstoffmangels} im
Vordergrund (Hypoxie): Der Patient leidet an Atemnot,
allerdings \E{ohne} Zyanose, begleitet von
Allgemeinsymptomen wie Schwindel, Kopfschmerz, Übelkeit und
Erbrechen, evtl. besteht auch eine Tachykardie.
Es können sich rasch Bewusstseinsstörungen bis hin zur
Bewusstlosigkeit entwickeln, auch Krämpfe sind gut möglich.
}{
\begin{itemize}
  \item Dyspnoe  \E{ohne}  Zyanose
  \item Tachykardie
  \item Schwindel, Kopfschmerz, Übelkeit, Erbrechen
  \item Bewusstseinsstörungen, Koma
  \item Krämpfe
  \item Pulsoxymeter zeigt falsche Werte an!
\end{itemize}
}{}{}{}


% M %%%%%%%%%%%%%%%%%%%%%%%%%%%%%%%%%%%%%%%%%%%%%%%%%%% M %
% Vergiftung mit Kohlenmonoxid
\ProcedurePrintDefault{MT58001C}


% \begin{figure}[H]
%     \begin{WidePar}
% \Parallel{
% \includegraphics[width=\linewidth]{./images/2006-01-31-gaswienheute2.png}
% }{
% \Captionof{figure}{\AuthNotesPicLegalOk{Gas Wien Heute
% Homepage}{OK, wenn richtig titriert}\ce{CO}-Vergiftung.
% Regelmäßig kommt es zu Unfällen durch defekte Gasthermen
% oder Heizlüfter. Bei diesem Einsatz war ein RTW des ASB
% Floridsdorf-Donaustadt beteiligt. Die dritte Patientin
% wurde erst \EE{nachdem die Feuerwehr \E{sämtliche} (!)
% Wohnungen des Wohnhauses aufgebrochen hatte} in der
% \E{darüberliegenden Wohnung} gefunden.
% \SourcePicture{Wien heute \url{http://wien.orf.at}}}
% }
% \end{WidePar}
% 
% \end{figure}

\Layout{27}
{}
{}
{}
{./images/2006-01-31-gaswienheute2.png}
{\EE{\ce{CO}-Vergiftung}:
Regel\-mäß\-ig kommt es zu Unfällen durch defekte Gasthermen
oder Heizlüfter. Bei diesem Einsatz war ein RTW des ASB
Floridsdorf-Donaustadt beteiligt. Die dritte Patientin
wurde erst \EE{nachdem die Feuerwehr \E{sämtliche}\,(!)
Wohnungen des Wohnhauses aufgebrochen hatte} in der
\E{darüberliegenden Wohnung} gefunden.}
{Wien heute, \url{http://wien.orf.at}}

\subsubsection{Kohlendioxid (\texorpdfstring{\ce{CO2}}{CO2}) -- Tod im Weinkeller}


Kohlendioxid ist ein farb- und geruchsloses Gas. Es ist
schwerer als Luft und sammelt sich z.\,B. in Mulden,
Weinkellern, Silos, etc. Es fällt als Gär- oder
Verbrennungsgas gehäuft in der Landwirtschaft und der
Industrie an.


\TextSlide{\TxSymptome}
{%
\begin{itemize}
  \item Kopfschmerz, Schwindel, Übelkeit
  \item Krämpfe
  \item Zyanose
  \item Hyperventilation, da der hohe
  \ce{CO2}-Spiegel im Blut das Atemzentrum stimuliert
\end{itemize}
}{
\begin{itemize}
  \item Kopfschmerz, Schwindel, Übelkeit
  \item Krämpfe
  \item Zyanose
  \item Hyperventilation
\end{itemize}
}{}{}{}


\MassnahmenMK{Spezielle
Maßnahmen}{Kohlendioxid-Vergiftung}{m:intoxikation-kohlendioxid}{
Allgemeine Maßnahmen bei Stickgas-Vergiftungen
(\prettyref{MK-m:am-stickgase})
}{}
% \MassnahmenNG{Spezielle
% Maßnahmen}{Kohlendioxid-Vergiftung}{\label{m:intoxikation-kohlendioxid}}{
% Allgemeine Maßnahmen bei Stickgas-Vergiftungen
% (\prettyref{m:am-stickgase})
% }

\Cave{
Bei Gasunfällen hat der \EE{Selbstschutz} Vorrang.
Oft ist die Bergung nur mit \EE{schwerem
Atemschutz} möglich ${\rightarrow}$ FEUERWEHR!}

\Fazit{Ein toter Helfer ist keine Hilfe.}

\subsection{Reizgase}

\Text{Typen}
{%
\begin{description}
  \item[Soforttyp] (z.\,B. Tränengas,...) Augenrötung, starker Hustenreiz
  \item[Verzögerungstyp] \Z{(z.\,B. Nitrosegase)} kaum
  Hustenreiz, Latenzphase (kaum Symptome),
  \EE{\VonBis{3}{24}\Hour* später: toxisches Lungenödem};
  bei Nichtbehandlung evtl. irreversible Lungenschäden
\end{description}
}{

}{}{}{}

\Layout{46}{Vorkommen}
{}{}{%
\begin{tabulary}{\linewidth}{LL}\addlinespace\toprule\addlinespace
\noindent\TabHeading{Vorkommen} & \noindent\TabHeading{Art}
\\\addlinespace\midrule\addlinespace
Chemische Industrie 
&
Nitrosegase,  Schwefelwasserstoffe, Ammoniak
\\\addlinespace
Kunststoffindustrie 
&
Chlorwasserstoffe,  Formaldehyd
\\\addlinespace
Düngemittelherstellung 
&
Ammoniak,  Nitrosegase, Schwefelwasserstoffe
\\\addlinespace
Haushalt 
&
Chlorgas (Kombination von best.  Haushaltsreinigern mit
Essig), Lacke,  Imprägniersprays, ...
\\\addlinespace\bottomrule
\end{tabulary}
}{Vorkommen von Reizgasen.}{}



% M %%%%%%%%%%%%%%%%%%%%%%%%%%%%%%%%%%%%%%%%%%%%%%%%%%% M %
% Vergiftung mit Reizgasen
\ProcedurePrintDefault{MT59XX2C}




\subsection{Kampfgase}

\Layout{60}{Abstand halten!}
{%
\noindent Kampfgase übersteigen die Expertise von zivilen Gesundheitseinrichtungen. 
Daher kurz und bündig: Hier müssen Spezialisten beigezogen werden.

Wichtig ist, auf Selbst- und Fremdschutz zu achten 
(Absperrmaßnahmen! Abstand halten! Wind beachten!), vgl. auch
\FullPrettyRef{topic:gefahrengutunfall}.
}{
\begin{itemize}
  \item Selbst- und Fremdschutz
  \item Spezialisten
  \item \FullPrettyRef{topic:gefahrengutunfall}
\end{itemize}
}
{./images/gabriel-sebastian-aass/schild-atemschutzmaske-abnehmen.jpg}
{}
{}


\clearpage
\section{Wovon man nicht trinken sollte \ldots}

\subsection{Einnahme von Säuren und Laugen}
\label{topic:veraetzung-intoxikation}

Verletzungen der Haut und Augen durch Säuren und
Laugen werden unter \FullPrettyRef{topic:veraetzung-trauma}
besprochen.

\Text
{\TxBeschreibung}
{%
\DefinitionInPara{Eine \T{Verätzung} ist eine lokale Schädigung der
Haut bzw. Schleimhaut aufgrund einer irreversiblen Zerstörung (Denaturierung) von
Eiweißstoffen. \Z{\Lang{Lat.} \I{Cauterisatio} (\I{Cauteris.})}}}
{}
{}
{}
{}
% 
% Betroffene Organe
% 
% \begin{itemize}
%     \item Magen-Darm-Trakt
%     \item Haut
%     \item Augen
% \end{itemize}

\TextSlide{\TxSymptome}
{
\begin{itemize}
    \item Ätzspuren (Mund/Rachen), Abrinnspuren
    \item Würgen
    \item Erbrechen
    \item Schmerzen
\end{itemize}
}{%
\begin{itemize}
  \item Ätz-, Abrinnspuren
  \item Würgen
  \item Erbrechen
  \item Schmerzen
\end{itemize}
}{}{}{}

\TextSlide{\TxGefahren}
{%
\begin{itemize}
  \item Ruptur/Perforation von Speiseröhre oder Magen
  \item Blutung
\end{itemize}
}{
\begin{itemize}
  \item Ruptur/Perforation Speiseröhre oder Magen
  \item Blutung
\end{itemize}
}{}{}{}



% M %%%%%%%%%%%%%%%%%%%%%%%%%%%%%%%%%%%%%%%%%%%%%%%%%%% M %
% Einnahme von Säuren oder Laugen
\ProcedurePrintDefault{MT54091C}


% Verätzung der Haut \prettyref{topic:veraetzung}



\subsection{Schaumbildner (Wasch-/Putzmittel)}


\TextSlide{{\TxBeschreibung}, {\TxSymptome} und {\TxGefahren}}{%
Schaumbildner vermindern die Oberflächenspannung, dadurch kommt es
zur Bildung von Schaum. Werden Schaumbildner
eingenommen, besteht durch die Schaumbildung eine \EE{hohe
Aspirationsgefahr} bzw. eine hohe Gefahr der \EE{Verlegung}
der Atemwege.
Zusätzlich kann der Patient über Übelkeit klagen bzw. erbrechen.

\bigskip

\bigskip

\bigskip
}{
\begin{itemize}
  \item Schaumbildung
  \begin{itemize}
    \item Aspirationsgefahr
    \item Verlegung der Atemwege
  \end{itemize}
  \item Übelkeit, Erbrechen
\end{itemize}
}{}{}{}


% M %%%%%%%%%%%%%%%%%%%%%%%%%%%%%%%%%%%%%%%%%%%%%%%%%%% M %
% Einnahme von schaumbildenden Substanzen
\ProcedurePrintDefault{MT17023C}



% 
% 
% \subsection{\Z{Reserviert}\A{Vergiftungen mit
% Insektiziden}}
% % TODO Vergiftungen mit Insektiziden
% 
% \subsection{\Z{Reserviert}\A{Vergiftungen durch giftige
% Tiere}}
% % TODO Vergiftungen durch giftige Tiere
% 
% 
% 

\ChapterEnd